\clearpage\pagestyle{matSimple}\phantomsection\label{guide}\pdfbookmark[0]{Guide de l’adulte}{bm-guide}
\matPageTitre{Guide de l’adulte}{Objectifs, manipulations, mots à faire dire et réponses}
{\fontsize{9.5}{13.5}\selectfont Pour chaque activité : l’objectif tiré du programme, ce que l’on peut faire avec de vrais objets avant la fiche, les mots que l’enfant est invité à dire, puis la réponse attendue. Les images en vert montrent une réponse correcte ; d’autres réponses peuvent être justes lorsque la consigne le permet.\par}
\par\needspace{62mm}\vspace{6mm}{\otSemi\fontsize{9}{11}\selectfont\color{matAccent}\otCaps[8]{Partie 1 · Compter jusqu’à dix}}\par\vspace{1mm}
\guideEntree{act-gs-1}{1}{Sept partout}{Reconnaître des collections de sept objets de tailles, de natures et de dispositions différentes.}{Faire des collections de sept objets très différents (sept perles, sept livres) et les comparer en les alignant.}{« Les ballons sont petits ou gros, rangés ou en vrac, mais il y en a toujours sept. »}{Les boîtes des étoiles, des ballons, des jetons et des fleurs. Il y a six pommes et huit poissons.}{gs-01}
\guideEntree{act-gs-2}{2}{Le bon nombre}{Dénombrer une collection jusqu’à dix et l’associer à son écriture chiffrée.}{Jeu du marchand : l’enfant compte des objets et pose la carte chiffrée qui convient.}{« Il y a dix étoiles : cinq et cinq. Je relie au 10. »}{Huit fraises avec 8, dix étoiles avec 10, six poissons avec 6, neuf bonbons avec 9.}{gs-02}
\guideEntree{act-gs-3}{3}{Deux groupes de points}{Dénombrer en utilisant des compositions : « quatre et quatre, huit ».}{Montrer rapidement deux constellations de dé ; l’enfant dit le total sans compter un à un, puis vérifie.}{« Je vois quatre et encore quatre : ça fait huit. Je sais que quatre et quatre font huit. »}{Huit, huit, dix et neuf. Pour la dernière carte, on peut garder six dans sa tête et surcompter : sept, huit, neuf.}{gs-03}
\guideEntree{act-gs-4}{4}{Tous les bonbons}{Parcourir une grande collection en marquant les éléments déjà comptés.}{Compter des jetons posés en vrac en les poussant d’un côté de la table au fur et à mesure.}{« Celui-là, je l’ai déjà compté : il a sa croix. Un, deux, trois… quatorze. »}{Quatorze bonbons. Chaque bonbon porte une seule croix ; aucun n’est oublié.}{gs-04}
\guideEntree{act-gs-5}{5}{Autant que le nombre}{Constituer une collection dont la quantité est donnée par son écriture chiffrée.}{Tirer une carte chiffrée et aller chercher exactement autant d’objets, en un seul voyage.}{« C’est un sept : je dessine sept ronds. Cinq et encore deux. »}{Sept, neuf et six ronds.}{gs-05}
\guideEntree{act-gs-6}{6}{Le plus grand nombre}{Comparer des quantités données par leur écriture chiffrée.}{Jeu de la bataille avec des cartes de un à dix : on dit qui a le plus grand nombre et pourquoi.}{« Neuf est plus grand que sept : sur la bande, il est plus loin. »}{9, 10, 8 et 6.}{gs-06}
\guideEntree{act-gs-7}{7}{Autant que le modèle}{Comparer des collections en les dénombrant, quand elles sont organisées différemment.}{Deux enfants ont des collections rangées différemment : ont-ils autant d’objets ? On compte pour le savoir.}{« Le modèle a huit pommes. Cette boîte aussi : huit. Il y en a autant. »}{La boîte du milieu (huit pommes). Les autres en ont sept et neuf.}{gs-07}
\par\needspace{62mm}\vspace{6mm}{\otSemi\fontsize{9}{11}\selectfont\color{matAccent}\otCaps[8]{Partie 2 · Composer les nombres}}\par\vspace{1mm}
\guideEntree{act-gs-8}{8}{Un de moins}{Réaliser une collection qui a un objet de moins et nommer le nombre avant et après le retrait.}{Faire une tour de sept cubes, en retirer un : « sept, j’en enlève un, six ».}{« Il y avait sept étoiles, j’en enlève une : il en reste six. Six vient juste avant sept. »}{6, 9 et 8 : quand on enlève un, on obtient le nombre qui vient juste avant.}{gs-08}
\guideEntree{act-gs-9}{9}{Juste avant, juste après}{Comprendre que chaque nombre s’obtient en ajoutant un au précédent.}{Réciter la comptine et s’arrêter : « Quel nombre vient juste après six ? Juste avant ? »}{« Avant six, il y a cinq. Après six, il y a sept : six et un, sept. »}{5 et 7 ; 8 et 10 ; 2 et 4 ; 1 et 3.}{gs-09}
\guideEntree{act-gs-10}{10}{Je garde dans ma tête}{Surcompter à partir d’un nombre donné.}{Cacher cinq jetons sous un gobelet, en ajouter trois à côté : « Cinq sont cachés. Combien en tout ? »}{« Je mets cinq dans ma tête : six, sept, huit. Il y en a huit. »}{8, 9 et 6. On part du nombre du dé et on continue la comptine, un mot par jeton.}{gs-10}
\guideEntree{act-gs-11}{11}{Les doubles des papillons}{Connaître les doubles : deux et deux, trois et trois, quatre et quatre, cinq et cinq.}{Plier une feuille peinte pour faire des papillons symétriques ; compter les taches d’une aile, puis des deux.}{« Trois et trois font six. Deux fois trois, six. Cinq et cinq font dix. »}{4, 6, 8 et 10.}{gs-11}
\guideEntree{act-gs-12}{12}{Les doigts des deux mains}{Composer les nombres jusqu’à dix avec les doigts des deux mains.}{Montrer huit avec ses doigts de plusieurs façons : cinq et trois, quatre et quatre…}{« Une main pleine, c’est cinq ; et encore trois : six, sept, huit. »}{8, 6, 10 et 9.}{gs-12}
\guideEntree{act-gs-13}{13}{La maison du huit}{Décomposer huit de plusieurs façons.}{Cacher une partie de huit jetons sous sa main : combien sont cachés ?}{« Huit, c’est cinq et trois. Huit, c’est six et deux. »}{Trois, cinq, deux et quatre points.}{gs-13}
\guideEntree{act-gs-14}{14}{Les dominos du sept}{Reconnaître différentes décompositions de sept.}{Trier un jeu de dominos : chercher tous ceux qui font sept.}{« Quatre et trois, sept. Six et un, sept aussi. Cinq et un, ça fait six. »}{Les dominos quatre-trois, six-un et cinq-deux.}{gs-14}
\par\needspace{62mm}\vspace{6mm}{\otSemi\fontsize{9}{11}\selectfont\color{matAccent}\otCaps[8]{Partie 3 · Écrire et réciter}}\par\vspace{1mm}
\guideEntree{act-gs-15}{15}{J’écris 7, 8, 9}{Écrire en chiffres les nombres de un à dix.}{Tracer le chiffre en grand dans l’air, puis sur une ardoise, en partant toujours du haut.}{« Le huit : je pars en haut, je fais un S, et je remonte. »}{Les chiffres sont tracés en partant du haut. Le sept peut s’écrire avec ou sans barre.}{gs-15}
\guideEntree{act-gs-16}{16}{J’écris jusqu’à 10}{Écrire en chiffres les nombres de un à dix pour dire une quantité.}{Écrire le nombre d’élèves présents, de jours de la semaine, dans des situations réelles de la classe.}{« Dix, ça s’écrit avec un 1 et un rond à côté. »}{Dix, sept et huit.}{gs-16}
\guideEntree{act-gs-17}{17}{Relie de 1 à 20}{Connaître la comptine numérique au-delà de dix et suivre la suite des nombres.}{Réciter la comptine jusqu’à trente en se déplaçant sur une piste numérotée dans la cour.}{« …dix-huit, dix-neuf, vingt. C’est un poisson ! »}{Un poisson. L’adulte aide à lire les nombres au-delà de dix si besoin : la comptine orale est l’objectif.}{gs-17}
\guideEntree{act-gs-18}{18}{Le compte à rebours}{Réciter la comptine numérique à rebours de dix à un.}{Faire décoller une fusée en carton : toute la classe compte à rebours, de dix jusqu’à un.}{« Dix, neuf, huit, sept, six, cinq, quatre, trois, deux, un… décollage ! »}{Il manque 8, 6, 3 et 1.}{gs-18}
\guideEntree{act-gs-19}{19}{De deux en deux}{Réciter la comptine de deux en deux, en partant de deux.}{Compter les élèves rangés par deux, les chaussures, les mains : « deux, quatre, six… ».}{« Deux, quatre, six, huit, dix ! » Puis jusqu’à vingt.}{La cinquième paire (dans la première rangée, la dernière). Il y a vingt chaussettes en tout.}{gs-19}
\guideEntree{act-gs-20}{20}{Le train part de douze}{Réciter la comptine en partant d’un nombre autre que un.}{Jeu du furet : un enfant dit un nombre, le suivant continue la comptine à partir de là.}{« Douze, treize, quatorze, quinze, seize. »}{Le cinquième wagon : douze, treize, quatorze, quinze, seize.}{gs-20}
\par\needspace{62mm}\vspace{6mm}{\otSemi\fontsize{9}{11}\selectfont\color{matAccent}\otCaps[8]{Partie 4 · Rangs et bande}}\par\vspace{1mm}
\guideEntree{act-gs-21}{21}{Dix cartes}{Déterminer un rang dans une suite de dix objets quand on change le point de départ.}{Aligner dix cartes entre un disque bleu et un disque rouge ; cacher un dessin sous une carte et faire décrire sa place.}{« Je pars du bleu : un, deux… sept. Je pars du rouge : un, deux, trois. »}{La septième carte en partant du bleu est coloriée ; la troisième en partant du rouge (la huitième en partant du bleu) est entourée.}{gs-21}
\guideEntree{act-gs-22}{22}{Les canards vont à la mare}{Repérer le premier, le deuxième, l’avant-dernier et le dernier d’une file.}{Faire une file : qui est le premier, le dernier ? Et si la file se retourne ?}{« Le premier est le plus près de la mare. Le dernier est tout au bout. »}{Le canard le plus proche de la mare en rouge, le suivant en jaune ; le dernier (le plus loin) en bleu et celui juste devant lui en vert.}{gs-22}
\guideEntree{act-gs-23}{23}{À la même place}{Repérer la position d’un objet dans une suite ordonnée pour la reproduire.}{Un enfant regarde un train modèle et dicte à un camarade où placer chaque symbole.}{« L’étoile est dans le deuxième wagon après la locomotive. Le rond dans le cinquième. »}{L’étoile dans le deuxième wagon, le rond dans le cinquième, le triangle dans le septième.}{gs-23}
\guideEntree{act-gs-24}{24}{La bande jusqu’à dix}{Compléter une bande numérique lacunaire jusqu’à dix.}{Fabriquer la bande numérique de la classe : une case par nombre, avec la quantité et le chiffre.}{« Après deux, il y a trois. Avant sept, il y a six. »}{3, 6, 9 ; 2, 5, 7, 10 ; 1, 4, 8.}{gs-24}
\guideEntree{act-gs-25}{25}{Sur la bande}{Placer différentes représentations des nombres jusqu’à dix sur la bande numérique.}{Accrocher sous la bande de la classe des cartes : dés, doigts, tours de cubes, collections.}{« La tour a sept cubes : je vais au 7. »}{Le dé six au 6, les huit doigts au 8, la tour de sept cubes au 7, les dix points au 10, le dé quatre au 4.}{gs-25}
\guideEntree{act-gs-26}{26}{Avancer, reculer}{Déterminer la position d’arrivée après un déplacement sur la bande numérique.}{Jeu de piste sur une grande bande au sol : avancer ou reculer en comptant les cases.}{« Je suis sur le quatre, j’avance de trois : cinq, six, sept. J’arrive sur le sept. »}{Le 7, le 6 et le 8. Avancer de trois revient à ajouter trois.}{gs-26}
\par\needspace{62mm}\vspace{6mm}{\otSemi\fontsize{9}{11}\selectfont\color{matAccent}\otCaps[8]{Partie 5 · Petits problèmes}}\par\vspace{1mm}
\guideEntree{act-gs-27}{27}{Les cubes dans la boîte}{Déterminer le tout dans un problème à trois parties.}{Mettre successivement des cubes dans une boîte opaque devant l’enfant, puis demander combien elle en contient.}{« Trois, et un, quatre ; et deux, cinq, six. Il y a six cubes. »}{Six cubes.}{gs-27}
\guideEntree{act-gs-28}{28}{Les oiseaux s’envolent}{Résoudre un problème de retrait en s’aidant d’une représentation.}{Accrocher sept pinces à linge sur une ficelle, en retirer trois et compter celles qui restent.}{« Sept, j’en enlève trois : il en reste quatre. Je le sais aussi : sept, c’est quatre et trois. »}{Trois oiseaux barrés (n’importe lesquels) ; il en reste quatre.}{gs-28}
\guideEntree{act-gs-29}{29}{Une carotte pour chacun}{Trouver la quantité à ajouter pour passer d’une collection à une autre.}{Huit enfants veulent un vélo, deux sont sortis : combien faut-il en sortir encore ?}{« Il y a huit lapins et deux carottes : il manque six carottes. Deux et six, huit. »}{Six carottes.}{gs-29}
\guideEntree{act-gs-30}{30}{Trois billes de plus}{Trouver une quantité à partir d’une autre et de l’écart entre les deux.}{Donner quatre jetons à un enfant ; un autre doit en avoir trois de plus : combien lui en donner ?}{« Le lapin a autant de billes que l’ours, et encore trois : quatre, cinq, six, sept. »}{Sept billes.}{gs-30}
\guideEntree{act-gs-31}{31}{Trois billes de moins}{Résoudre un problème d’écart dont l’énoncé dit « de moins » alors qu’il faut ajouter.}{Jouer la scène avec des jetons : si le chat en a trois de moins, c’est que la souris en a trois de plus.}{« Le chat en a moins. Donc la souris en a plus : cinq, et encore trois, huit. »}{Huit billes. Attention : le mot « moins » ne veut pas dire qu’il faut enlever.}{gs-31}
\guideEntree{act-gs-32}{32}{D’où est parti le pion ?}{Trouver la position de départ à partir de la position d’arrivée et du déplacement.}{Sur une piste au sol : « Je suis arrivé ici en avançant de trois. D’où suis-je parti ? »}{« Il a avancé de trois pour arriver sur neuf. Je recule de trois : huit, sept, six. Il était sur six. »}{Il est parti de 6 (il a avancé de trois jusqu’à 9), puis de 7 (il a reculé de deux jusqu’à 5).}{gs-32}
\guideEntree{act-gs-33}{33}{Le goûter}{Déterminer le tout dans un problème de groupements.}{Quatre assiettes sur une table, des jetons sur une autre : aller chercher en un seul voyage de quoi mettre deux gâteaux par assiette.}{« Deux, quatre, six, huit. Il faut huit gâteaux. »}{Huit gâteaux : deux, quatre, six, huit.}{gs-33}
\guideEntree{act-gs-34}{34}{Partager sept gâteaux}{Partager équitablement avec un reste.}{Distribuer sept jetons à deux poupées, un à un : il en reste un qu’on ne peut pas donner sans faire de jaloux.}{« Un pour toi, un pour toi… Chacun en a trois, et il en reste un. »}{Trois gâteaux pour chaque ours et un gâteau qui reste.}{gs-34}
\guideEntree{act-gs-35}{35}{Les enveloppes}{Trouver le nombre de parts dans un problème de groupements.}{Ranger dix images dans des enveloppes, deux par enveloppe, et compter les enveloppes utilisées.}{« Deux, quatre, six, huit, dix : cinq enveloppes. »}{Cinq enveloppes. On peut regrouper les images de bien des façons, par paires.}{gs-35}
\par\needspace{62mm}\vspace{6mm}{\otSemi\fontsize{9}{11}\selectfont\color{matAccent}\otCaps[8]{Partie 6 · Formes et solides}}\par\vspace{1mm}
\guideEntree{act-gs-36}{36}{Quatre familles de formes}{Reconnaître et classer carrés, rectangles, triangles et disques, quelles que soient leur taille et leur orientation.}{Trier des formes en carton dans quatre boîtes ; les décrire : nombre de côtés, côtés pareils ou non.}{« Le carré a quatre côtés tous pareils. Le rectangle a quatre côtés : deux longs et deux courts. »}{Quatre formes de chaque famille. Les formes penchées gardent leur nom.}{gs-36}
\guideEntree{act-gs-37}{37}{Combien de côtés ?}{Décrire des formes planes : compter leurs côtés.}{Faire le tour d’une forme en bois avec le doigt en comptant les côtés.}{« Le triangle a trois côtés. Le carré et le rectangle en ont quatre. »}{3, 4, 4 ; 3, 4, 4.}{gs-37}
\guideEntree{act-gs-38}{38}{Les faces du cube}{Décrire un cube : six faces carrées, toutes pareilles.}{Observer un vrai cube : poser une gommette sur chaque face en comptant. Comparer avec un pavé.}{« Le cube a six faces. Ce sont des carrés, tous de la même taille. »}{Les six petits carrés pareils. Le grand carré est trop grand ; les autres formes ne sont pas des carrés.}{gs-38}
\guideEntree{act-gs-39}{39}{Les faces de la pyramide}{Décrire une pyramide à base carrée : un carré et quatre triangles.}{Observer une vraie pyramide à base carrée, compter ses faces, les faire « tamponner » dans la peinture.}{« La pyramide a une face carrée et quatre faces triangulaires toutes pareilles. »}{Le grand carré et les quatre triangles pareils. Le petit carré, le petit triangle, le disque et le rectangle ne conviennent pas.}{gs-39}
\guideEntree{act-gs-40}{40}{Le tangram}{Reproduire un assemblage de sept formes, même s’il n’est pas à la même taille.}{Reconstituer un vrai tangram sur son modèle, puis sur un modèle plus grand ou plus petit.}{« Les deux grands triangles sont à gauche et en bas. Le carré est penché, au milieu. »}{Chaque pièce garde sa place et sa couleur ; seule la taille change.}{gs-40}
\guideEntree{act-gs-41}{41}{Avec la règle}{Utiliser la règle pour tracer des traits droits.}{Apprendre à tenir la règle : la main à plat au milieu, le crayon qui glisse le long du bord.}{« Je pose la règle sur les deux points, je la tiens bien et je trace. »}{Une maison à cinq côtés et une étoile à cinq branches.}{gs-41}
\par\needspace{62mm}\vspace{6mm}{\otSemi\fontsize{9}{11}\selectfont\color{matAccent}\otCaps[8]{Partie 7 · Longueurs et masses}}\par\vspace{1mm}
\guideEntree{act-gs-42}{42}{La bande témoin}{Comparer des longueurs avec une bande témoin quand on ne peut pas rapprocher les objets.}{Comparer la largeur de la porte et celle de la fenêtre avec une ficelle ou une bande de papier.}{« Je reporte la longueur du crayon bleu sur ma bande. Le rouge dépasse : il est plus long. »}{Le crayon rouge. Les écarts sont petits : il faut reporter les longueurs pour en être sûr.}{gs-42}
\guideEntree{act-gs-43}{43}{Du plus court au plus long}{Ranger cinq objets selon leur longueur.}{Ranger cinq bandes de papier en escalier, en les mettant à la même origine.}{« Le vert est le plus court : un. Le jaune est le plus long : cinq. »}{Dans l’ordre des rubans : 3, 1, 5, 2, 4.}{gs-43}
\guideEntree{act-gs-44}{44}{Aussi long}{Produire un trait de même longueur qu’un objet donné.}{Couper une ficelle aussi longue qu’un crayon, puis aussi longue que la table.}{« Je mets ma règle sous le crayon, je repère où il s’arrête, puis je trace jusque-là. »}{Deux traits qui commencent au point et finissent juste sous la fin du crayon et du ruban.}{gs-44}
\guideEntree{act-gs-45}{45}{Le plus lourd des trois}{Ordonner les masses de trois objets à partir de deux comparaisons.}{Comparer trois boîtes de même taille sur une balance de Roberval, deux par deux.}{« Le rouge est plus lourd que le bleu. Le bleu est plus lourd que le jaune. Donc le rouge est le plus lourd. »}{Le paquet rouge est le plus lourd, le jaune le plus léger.}{gs-45}
\guideEntree{act-gs-46}{46}{L’équilibre}{Réaliser l’équilibre sur une balance de Roberval.}{Avec une vraie balance et des cubes identiques, chercher à rendre la balance droite.}{« À gauche, il y a quatre cubes ; à droite, deux. Il faut en ajouter deux pour avoir autant. »}{2, puis 3.}{gs-46}
\par\needspace{62mm}\vspace{6mm}{\otSemi\fontsize{9}{11}\selectfont\color{matAccent}\otCaps[8]{Partie 8 · Les motifs}}\par\vspace{1mm}
\guideEntree{act-gs-47}{47}{Ça grandit}{Repérer et prolonger un motif évolutif.}{Faire des rythmes qui grandissent : un tape, un claque ; deux tapes, deux claques…}{« Un carré, un rond ; deux carrés, deux ronds ; trois carrés, trois ronds… puis quatre carrés. »}{Première rangée : quatre carrés. Deuxième rangée : une étoile, puis quatre ronds.}{gs-47}
\guideEntree{act-gs-48}{48}{Plus grand, plus petit}{Prolonger un motif évolutif et dire sa règle.}{Construire des tours de cubes de plus en plus hautes, puis de plus en plus basses.}{« Un cube de plus à chaque fois : la suivante a cinq cubes. En bas, une étoile de moins : il en reste deux. »}{Une tour de cinq cubes ; une pile de deux étoiles.}{gs-48}
\guideEntree{act-gs-49}{49}{Le même motif, autrement}{Transcrire un motif avec d’autres symboles en gardant sa structure.}{Traduire un collier rouge-bleu en sons (tape, claque) ou en gestes (debout, assis).}{« Carré, triangle, carré, triangle… devient étoile, rond, étoile, rond. »}{Première rangée : étoile, rond, étoile, rond, étoile, rond. Deuxième : triangle, triangle, étoile, trois fois.}{gs-49}
\guideEntree{act-gs-50}{50}{La même règle}{Reconnaître des motifs qui ont la même structure avec des éléments différents.}{Comparer « tape, cuisse, cuisse » et un collier « rouge, bleu, bleu » : c’est la même règle.}{« Un, puis deux pareils ; un, puis deux pareils. »}{La première (étoile, rond, rond) et la troisième (triangle, carré, carré) : une forme, puis deux formes pareilles.}{gs-50}
\guideEntree{act-gs-51}{51}{J’invente mon collier}{Créer un motif et le décrire pour qu’un autre puisse le reproduire.}{Un enfant invente un collier et le décrit à un camarade qui le fabrique sans le voir.}{« Ma règle : deux rouges, une bleue, et on recommence. »}{Toutes les réponses sont justes si les couleurs suivent la règle annoncée. L’image montre un exemple.}{gs-51}
\par\needspace{62mm}\vspace{6mm}{\otSemi\fontsize{9}{11}\selectfont\color{matAccent}\otCaps[8]{Partie 9 · Se repérer}}\par\vspace{1mm}
\guideEntree{act-gs-52}{52}{Vers la droite}{Distinguer la droite et la gauche.}{Jeu de Jacques a dit : « Lève la main droite, tourne vers la gauche… »}{« Ce canard regarde vers la droite, comme la flèche. Ce poisson regarde vers la gauche. »}{Le premier canard, l’oiseau, l’éléphant, la souris et le poisson orange regardent vers la droite.}{gs-52}
\guideEntree{act-gs-53}{53}{À gauche, à droite}{Situer un objet par rapport à un autre avec « à gauche de », « à droite de ».}{Ranger des jouets sur une étagère en suivant des consignes : « Mets la voiture à gauche de l’ours. »}{« La pomme est à droite du chat. La voiture est à gauche du lapin. »}{La pomme est entourée, la voiture barrée.}{gs-53}
\guideEntree{act-gs-54}{54}{Le code secret}{Réaliser un déplacement à partir de son codage par des flèches.}{Dans la cour, sur un quadrillage tracé à la craie, avancer d’une case par flèche lue par un camarade.}{« Droite, droite, haut, haut, droite, droite, haut, gauche. »}{On arrive dans la case juste sous l’arbre.}{gs-54}
\guideEntree{act-gs-55}{55}{J’écris le chemin}{Coder un déplacement à l’aide de flèches.}{Un enfant se déplace sur un quadrillage au sol ; un autre code son trajet avec des cartes-flèches.}{« La souris va à droite, à droite, puis elle descend… »}{Droite, droite, bas, droite, bas, bas, droite, droite, droite, bas.}{gs-55}
\guideEntree{act-gs-56}{56}{Sur le quadrillage}{Repérer des cases sur un quadrillage pour reproduire un dessin.}{Reproduire un dessin fait de gommettes sur une grille, en comptant les cases depuis un coin.}{« Le toit rouge est en haut, au milieu. La porte marron est sur la deuxième ligne. »}{Une maison : un toit rouge, des murs jaunes, une porte marron, et une bande d’herbe verte en bas.}{gs-56}
